% GENERATED FILE. Edit canonical lessons, then run npm run generate:books.
% canonical-source-hash: fnv1a64:6799f6b5e8532df0
% canonical-lessons: GE-C426-rezeption, GE-C426-anreisen, GE-C426-abreisen, GE-C426-aufenthalt, GE-C426-auschecken

\chapter{An der Rezeption}
\label{ch:ge-an-der-rezeption}
\hlchaptermodality{426}

\begin{chapteropening}
\textbf{By the end of this chapter:} \emph{I can arrive at and leave a hotel: the reception, arriving, leaving, the stay and checking out.}

Complete the last lesson of chapter 426: I can arrive at and leave a hotel: the reception, arriving, leaving, the stay and checking out.
\end{chapteropening}

\section[die Rezeption]{die Rezeption — the reception (desk)}

\label{lesson:GE-C426-rezeption}



\begin{quote}
Before the new one: say the German for to put up, then the German for to take down.
\end{quote}

\subsection*{You'll want to know: die Rezeption}
\textbf{die Rezeption} — \textquotedblleft{}the reception (desk)\textquotedblright{}.

\begin{quote}
Den Schlüssel bekommt man an der Rezeption.
\end{quote}

\begin{quote}
\emph{You get the key at reception.}
\end{quote}

Plural: \emph{die Rezeptionen}. Also \emph{der Empfang}.

\begin{cousinweb}
\emph{Rezeption} is a learned borrowing from Latin \emph{receptiō} 'a receiving' (EtymDB), from \emph{recipere} 'to take back, to receive' (Lewis \& Short), probably through French \emph{réception}.
\end{cousinweb}

\begin{grammarlens}[title={What you've built}]
One more word for telling, explaining and describing at B1.
\end{grammarlens}

\subsection*{Your turn}
\begin{itemize}
\raggedright
  \item \emph{Say it:} \emph{die Rezeption}
  \item \emph{Say it:} \emph{die Rezeption}, once more
  \item \emph{Say it:} a sentence of your own with \emph{die Rezeption}
  \item \emph{Recall:} say the German for to put up, then the German for to take down, then say \emph{die Rezeption} again
\end{itemize}

\begin{tcolorbox}[breakable,colback=teal!4,colframe=teal!35!black,title={Before you move on}]
What does die Rezeption mean? (\textquotedblleft{}The reception (desk)\textquotedblright{}.) Say it once more.
\end{tcolorbox}
\section[anreisen]{anreisen — to arrive (at a place you are staying)}

\label{lesson:GE-C426-anreisen}



\begin{quote}
Before the new one: say the German for to take down, then the German for the reception.
\end{quote}

\subsection*{You'll want to know: anreisen}
\textbf{anreisen} — \textquotedblleft{}to arrive (at a place you are staying)\textquotedblright{}.

\begin{quote}
Die Gäste reisen am Freitag an und bleiben eine Woche.
\end{quote}

\begin{quote}
\emph{The guests arrive on Friday and stay for a week.}
\end{quote}

Separable; Perfekt with \emph{sein}. \emph{Der Anreisetag} — the day of arrival.

\begin{cousinweb}
\emph{anreisen} is built from \emph{an} 'at, to' and \emph{reisen} 'to travel', inherited from Old High German \emph{reisōn} 'to set out' (EtymDB).
\end{cousinweb}

\begin{grammarlens}[title={What you've built}]
One more word for telling, explaining and describing at B1.
\end{grammarlens}

\subsection*{Your turn}
\begin{itemize}
\raggedright
  \item \emph{Say it:} \emph{anreisen}
  \item \emph{Say it:} \emph{anreisen}, once more
  \item \emph{Say it:} a sentence of your own with \emph{anreisen}
  \item \emph{Recall:} say the German for to take down, then the German for the reception, then say \emph{anreisen} again
\end{itemize}

\begin{tcolorbox}[breakable,colback=teal!4,colframe=teal!35!black,title={Before you move on}]
What does anreisen mean? (\textquotedblleft{}To arrive (at a place you are staying)\textquotedblright{}.) Say it once more.
\end{tcolorbox}
\section[abreisen]{abreisen — to leave, to depart (at the end of a stay)}

\label{lesson:GE-C426-abreisen}



\begin{quote}
Before the new one: say the German for the reception, then the German for to arrive.
\end{quote}

\subsection*{You'll want to know: abreisen}
\textbf{abreisen} — \textquotedblleft{}to leave, to depart (at the end of a stay)\textquotedblright{}.

\begin{quote}
Wir reisen morgen früh ab.
\end{quote}

\begin{quote}
\emph{We are leaving early tomorrow.}
\end{quote}

Separable; Perfekt with \emph{sein}: \emph{sie sind schon abgereist}.

\begin{cousinweb}
\emph{abreisen} is built from \emph{ab} 'off, away' and \emph{reisen} 'to travel', inherited from Old High German \emph{reisōn} 'to set out' (EtymDB).
\end{cousinweb}

\begin{grammarlens}[title={What you've built}]
One more word for telling, explaining and describing at B1.
\end{grammarlens}

\subsection*{Your turn}
\begin{itemize}
\raggedright
  \item \emph{Say it:} \emph{abreisen}
  \item \emph{Say it:} \emph{abreisen}, once more
  \item \emph{Say it:} a sentence of your own with \emph{abreisen}
  \item \emph{Recall:} say the German for the reception, then the German for to arrive, then say \emph{abreisen} again
\end{itemize}

\begin{tcolorbox}[breakable,colback=teal!4,colframe=teal!35!black,title={Before you move on}]
What does abreisen mean? (\textquotedblleft{}To leave, to depart (at the end of a stay)\textquotedblright{}.) Say it once more.
\end{tcolorbox}
\section[der Aufenthalt]{der Aufenthalt — the stay}

\label{lesson:GE-C426-aufenthalt}



\begin{quote}
Before the new one: say the German for to arrive, then the German for to leave.
\end{quote}

\subsection*{You'll want to know: der Aufenthalt}
\textbf{der Aufenthalt} — \textquotedblleft{}the stay\textquotedblright{}.

\begin{quote}
Wir wünschen Ihnen einen schönen Aufenthalt.
\end{quote}

\begin{quote}
\emph{We wish you a pleasant stay.}
\end{quote}

Plural: \emph{die Aufenthalte}. At a station, \emph{der Aufenthalt} is also how long a train stops.

\begin{cousinweb}
\emph{Aufenthalt} is inherited from Middle High German \emph{ūfenthalt} 'stay' (EtymDB), built on \emph{halten} 'to hold, to stop', from Proto-Germanic \emph{haldaną}.
\end{cousinweb}

\begin{grammarlens}[title={What you've built}]
One more word for telling, explaining and describing at B1.
\end{grammarlens}

\subsection*{Your turn}
\begin{itemize}
\raggedright
  \item \emph{Say it:} \emph{der Aufenthalt}
  \item \emph{Say it:} \emph{der Aufenthalt}, once more
  \item \emph{Say it:} a sentence of your own with \emph{der Aufenthalt}
  \item \emph{Recall:} say the German for to arrive, then the German for to leave, then say \emph{der Aufenthalt} again
\end{itemize}

\begin{tcolorbox}[breakable,colback=teal!4,colframe=teal!35!black,title={Before you move on}]
What does der Aufenthalt mean? (\textquotedblleft{}The stay\textquotedblright{}.) Say it once more.
\end{tcolorbox}
\section[auschecken]{auschecken — to check out}

\label{lesson:GE-C426-auschecken}



\begin{quote}
Before the new one: say the German for to leave, then the German for the stay.
\end{quote}

\subsection*{You'll want to know: auschecken}
\textbf{auschecken} — \textquotedblleft{}to check out\textquotedblright{}.

\begin{quote}
Bis elf Uhr müssen die Gäste auschecken.
\end{quote}

\begin{quote}
\emph{The guests have to check out by eleven o'clock.}
\end{quote}

Separable: \emph{ich checke aus}, \emph{ich habe ausgecheckt}.

\begin{cousinweb}
\emph{auschecken} is built from \emph{aus} 'out' and \emph{checken}, borrowed from English \emph{check} (EtymDB), on the model of English \emph{check out}.
\end{cousinweb}

\begin{grammarlens}[title={What you've built}]
One more word for telling, explaining and describing at B1.
\end{grammarlens}

\subsection*{Your turn}
\begin{itemize}
\raggedright
  \item \emph{Say it:} \emph{auschecken}
  \item \emph{Say it:} \emph{auschecken}, once more
  \item \emph{Say it:} a sentence of your own with \emph{auschecken}
  \item \emph{Recall:} say the German for to leave, then the German for the stay, then say \emph{auschecken} again
\end{itemize}

\begin{tcolorbox}[breakable,colback=teal!4,colframe=teal!35!black,title={Before you move on}]
What does auschecken mean? (\textquotedblleft{}To check out\textquotedblright{}.) Say it once more.
\end{tcolorbox}
